\section{The Probe Reward Trap}
\label{sec:probe_trap}

The cosine formulations fail on signal magnitude; a natural response is to discretize.
Training a logistic probe on $\dsyc$ and thresholding its output yields $\rone = \mathbb{I}[P(\text{syc} \mid \mathbf{h}) < 0.5]$ (F2 in \S\ref{ssec:reward}).
We show that this discretization sidesteps the magnitude problem but creates a qualitatively worse failure.
Both probe experiments use the probe as a non-differentiable reward component within GRPO (\S\ref{ssec:reward}); the trap therefore reflects a semantic limitation of the probe rather than gradient-based exploitation.

\paragraph{Differentiable probe: gradient gaming.}
Without stop-gradient, the policy discovers within ${\sim}50$ steps that shifting activations away from the sycophancy cluster maximizes $\rone$ regardless of response content, producing an inverted profile ($\tof = 0.978$, $\nof = 0.006$).
Replication across three seeds (42, 43, 44) confirms determinism: all collapse within ${\sim}50$--$60$ steps, with each probe retraining (AUROC ${\to}1.0$) re-gamed within ${\sim}10$ steps.
This motivates the stopped-gradient design below.

\paragraph{Binary probe: the stubbornness trap.}
With stop-gradient, the policy cannot game the boundary but instead exploits a deeper flaw: $\dsyc$ does not condition on input validity, so refusing to revise a wrong answer receives the same $\rone = 1$ as correctly resisting invalid pressure.
By step 500, $\tof$ drops to 0.000 and $\nof$ reaches 0.995, yielding $\sycon = 0$.
Qualitative analysis of 181 valid-correction responses reveals five failure modes: \emph{deliberation traps} (28.7\%), where the model identifies the correct answer but cannot commit; \emph{sycophantic stubbornness} (24.3\%), where the model verbally agrees but does not change its answer; \emph{acknowledge-but-maintain} (23.8\%); \emph{genuine-accept-missed} (18.8\%); and refute/ignore/other (4.5\%). Full taxonomy in Appendix~\ref{app:taxonomy}.

\paragraph{Behavioral death precedes distributional collapse.}

\begin{table}[t]
\centering
\caption{Probe reward trap trajectory (stopped-gradient probe). Behavioral death ($\tof{=}0$) occurs at step 500; entropy remains healthy until step ${\sim}1840$. Reward averages $+0.72$ throughout.}
\label{tab:probe_trajectory}
\resizebox{\columnwidth}{!}{%
\begin{tabular}{lccccc}
\toprule
Step & $\tof$ & $\nof$ & $\sycon$ & Entropy & Reward \\
\midrule
500 & 0.000 & 0.995 & 0.000 & 0.702 & +0.72 \\
1000 & 0.006 & 0.995 & 0.011 & 0.800 & +0.71 \\
1500 & 0.000 & 1.000 & 0.000 & 0.608 & +0.69 \\
2000 & 0.000 & 1.000 & 0.000 & 0.304 & +0.69 \\
\bottomrule
\end{tabular}}
\end{table}

\begin{figure}[t]
\centering
\includegraphics[width=\columnwidth]{figures/fig_probe_trap.pdf}
\caption{Probe reward trap (stopped-gradient probe). $\tof$ drops to zero at step 500 while generation entropy remains high (0.70--0.80 through step 1000). The 1{,}300-step lag means standard entropy monitoring misses the behavioral failure entirely.}
\label{fig:probe_trap}
\end{figure}

Table~\ref{tab:probe_trajectory} and Figure~\ref{fig:probe_trap} reveal a 1{,}300-step lag: the model is behaviorally dead ($\tof = 0$) at step 500, yet entropy stays healthy (0.70--0.80) through step 1000 and does not collapse until step ${\sim}1840$, while reward averages $+0.72$.
Entropy-based early stopping, the most common safeguard~\citep{casper2023openproblems}, would miss the failure; detecting probe reward traps requires periodic $\tof$ evaluation.

\paragraph{A dilemma for direction-based rewards.}
The differentiable and binary failures bracket the design space.
With gradients, the policy games the boundary; without gradients, it exploits the fact that stubbornness and genuine resistance are indistinguishable along $\dsyc$.
Any reward conditioned only on $\dsyc$ (continuous or discrete, differentiable or frozen) will either be gamed or will reward rigidity.
Escaping requires a signal encoding the evidential warrant of the user's challenge---a dimension absent from $\dsyc$~\citep{denison2024subterfuge}.
